\chapter{Energy and Momentum of Rotating Systems}

\mathTopics{\pgRef{chap:vectors}}

\section{Rotational Kinetic Energy}

Similarly to the formula for linear kinetic energy in
\pgRef{eq:linear-kinetic-energy},
the formula for \textbf{rotational kinetic energy} is
\begin{equation}
    E_\mathrm{k} = \frac{1}{2} I \omega^2,
\end{equation}
where \(E_\mathrm{k}\) is the kinetic energy of the rotating object,
\(I\) is the moment of inertia of the object spinning around its axis
of rotation,
and \(\omega\) is the angular velocity.

\section{Rotational Work}

The \textbf{rotational work} \(W\) is given by
\begin{equation}
    W = \int \tau(\theta) \odif{\theta}.
\end{equation}

For example, let us find the the work done by gravity
as a rod swings from horizontal to vertical.
Let \(r\) be the distance between the axis of rotation and the rod's
center of mass
and \(m\) be the mass of the rod.
Then the work \(W\) is
\begin{equation}
    \begin{aligned}
        W &= \int_0^\frac{\pi}{2} \frac{1}{2} m g \cos\theta \odif{\theta} \\
        &= \frac{3}{2} g \ab\Big[\sin\theta]_0^\frac{\pi}{2} \\
        &= \frac{3}{2} g.
    \end{aligned}
\end{equation}

\section{Angular Momentum}

Unsurprisingly, \textbf{angular momentum} is the rotational analogue
to linear momentum.
It is given by the equation
\begin{equation}
    \vecb{L} = I \vecb{\omega}.
\end{equation}

Objects moving in a straight lines also have angular momentum.
The formula for such objects is
\begin{subequations}
    \begin{align}
        \vecb{L} &= \vecb{r} \times \vecb{p} \\
        L &= m v r \sin\theta. \label{eq:angular-momentum-linear-velocity}
    \end{align}
\end{subequations}

\subsection{Conservation of Angular Momentum}

Similarly to linear momentum, angular momentum must be conserved.

The most famous demonstration of \textbf{conservation of angular momentum}
is a person pulling in their arms while they are spinning.
Pulling in their arms decreases the radius between their arms and
center of mass,
reducing moment of inertia.
Since angular momentum must be conserved,
their angular velocity increases to compensate.

\section{Orbits} \label{sec:orbits}

\textbf{Kepler's laws of planetary motion}
describe the motion of planets orbiting around a star.
\begin{enumerate}
    \item Planetary orbits are elliptical,
        with the star at one of the ellipse's foci.
    \item \label{itm:keplers-2nd-law}
        The line between the star and an orbiting planet
        sweeps across equal amounts of area in equal amounts of time.
    \item The square of the orbital period \(T\) is proportional to
        the cube of the semi-major axis \(r\):
        \begin{equation}
            T^2 \propto r^3
        \end{equation}
        See \pgRef{subsec:orbital-velocity} for the derivation.
        (Strictly speaking, the derivation assumes a circular orbit,
            but the formula for the perimeter of an ellipse is rather messy,
        so let us pretend it is circular.)
\end{enumerate}

In elliptical \textbf{orbits}, both objects orbit around the center of mass,
in this context called the \textbf{barycenter}.

Imagine a planet of mass \(m\) orbitting around a star.
There are two points in its orbit: A and B.
Let
\begin{itemize}
    \item \(L_\mathrm{A}\) and \(L_\mathrm{B}\) be the angular
        momenta at points A and B, respectively.
    \item \(v_\mathrm{A}\) and \(v_\mathrm{B}\) be the velocities at
        points A and B, respectively.
    \item \(r_\mathrm{A}\) and \(r_\mathrm{B}\) be the distances from
        the star to points A and B, respectively.
        Assume \(r_\mathrm{A} < r_\mathrm{B}\).
\end{itemize}
\begin{figure}[ht]
    \centering
    \begin{tikzpicture}
    \tikzmath{
        \a = 4;
        \b = 3;
        \c = sqrt(\a^2 - \b^2);
    }

    \draw[-{>[scale=2]}] (0.0, \b) arc [
        start angle=90, end angle=460, x radius=\a, y radius=\b
    ];
    \fill (\c, 0.0) circle (0.25) node[above=3.0mm] {Star};

    \fill ({\a*cos(80)}, {\b*sin(80)}) circle (0.125) node[above=3.0mm] {\(m\)};
    \draw[dotted, line width = 0.25mm, line cap=round] (-\a, 0.0)
    circle (0.125) node[left=1.5mm] {B};
    \draw[dotted, line width = 0.25mm, line cap=round] (\a, 0.0)
    circle (0.125) node[right=1.5mm] {A};
\end{tikzpicture}

    \caption{
        (Not drawn to scale)
        Diagram of a planet with mass \(m\) orbitting around a star.
        Point A is closer to the star than point B.
    }
\end{figure}

From the conservation of angular momentum
and \pgRef{eq:angular-momentum-linear-velocity},
we know
\begin{align}
    L_\mathrm{A} &= L_\mathrm{B}, \\
    m v_\mathrm{A} r_\mathrm{A} &= m v_\mathrm{B} r_\mathrm{B}.
\end{align}
If we solve for either velocity we find
\begin{equation}
    \begin{aligned}
        v_\mathrm{A} &= v_\mathrm{B} \frac{r_\mathrm{B}}{r_\mathrm{A}}, \\
        v_\mathrm{B} &= v_\mathrm{A} \frac{r_\mathrm{A}}{r_\mathrm{B}}.
    \end{aligned}
\end{equation}
The above equations suggest that the planet moves faster at point A,
when it is closer to the star, than at point B,
when it is farther from the star.

This conclusion is actually a consequence of
\hyperref[itm:keplers-2nd-law]{Kepler's second law of planetary
motion} (\cpageref{itm:keplers-2nd-law}).
Since planetary motion is an ellipse,
for the same amount of time to sweep the same area,
the planet would have to travel faster to compensate for the smaller radius.

Mechanical energy of orbiting bodies:
\begin{gather}
    E = U + E_\mathrm{K} \\
    \begin{aligned}
        U &= - \int F \odif{r} \\
        &= - \int G \frac{m_1 m_2}{r^2} \odif{r} \\
        &= -G \frac{m_1 m_2}{r}
    \end{aligned}
\end{gather}
\begin{gather}
    \begin{aligned}
        E_\mathrm{K} = \frac{1}{2} m v^2 \nonumber \\
        a_\mathrm{c} &= \frac{v^2}{r} \\
        \therefore v^2 &= a_\mathrm{c} r
    \end{aligned}
    E_\mathrm{K} = \frac{1}{2} m a_\mathrm{c} r \\
    m a_\mathrm{c} = F_\mathrm{g} = G \frac{m_1 m_2}{r^2} \\
    \begin{aligned}
        E_\mathrm{K} &= \frac{1}{2} \frac{G m_1 m_2}{r^2} r \\
        &= \frac{G m_1 m_2}{2 r}
    \end{aligned}
\end{gather}
\begin{equation}
    \begin{aligned}
        E &= - \frac{G m_1 m_2}{r} + \frac{G m_1 m_2}{2 r} \\
        &= - \frac{G m_1 m_2}{2 r}
    \end{aligned}
\end{equation}

\subsection{Escape Velocities}

Let
\begin{itemize}
    \item \(M\) be the mass of the larger body, e.g., a planet.
    \item \(m\) be the mass of the escaping body, e.g., a spaceship.
    \item \(r\) the radius of the larger body.
    \item \(v\) be the escape velocity.
\end{itemize}
\(- \frac{G M m}{r} + \frac{1}{2} m v^2\) is the total energy.
To escape, you need both 0 gravitational potential energy and 0 kinetic energy.
Therefore, we find that
\begin{equation}
    \begin{aligned}
        - \frac{G M m}{r} + \frac{1}{2} m v^2 &= 0 \\
        \frac{1}{2} m v^2 &= \frac{G M m}{r} \\
        v &= \sqrt{\frac{2 G M}{r}}.
    \end{aligned}
\end{equation}
